
% ════════════════════════════════════════════
% 模型建立（5.4.3.模型建立.tex，三级子块）
% ════════════════════════════════════════════
\subsubsection{计及实时电价的两阶段决策模型}

问题四与问题二、问题三共用同一套模型结构，区别只在电价参数取逐间隔波动电价 $p_{d,k}$ 而非固定电价 $p_k$。

\textbf{步骤1：波动电价下的逐日计划模型}

对应问题二，第 $d$ 天的逐日计划购电模型的目标函数、功率平衡与储能状态递推约束为
\begin{equation}
\begin{aligned}
&\min \; Z_d=\sum_{k=0}^{143} g_{d,k}\,p_{d,k}\,\Delta t\\
&g_{d,k}+S_{d,k}+d_{d,k}-c_{d,k}\geq L_{d,k}\\
&E_{d,k+1}=E_{d,k}+\eta\,c_{d,k}\,\Delta t-\frac{1}{\eta}\,d_{d,k}\,\Delta t
\end{aligned}
\end{equation}
容量与功率约束、日间储电量传递关系与首日初值（同问题二）为
\begin{equation}
\begin{aligned}
&E_{\min}\leq E_{d,k}\leq E_{\max},\quad 0\leq c_{d,k},d_{d,k}\leq P_{\max},\quad g_{d,k}\geq 0\\
&E_{d+1,0}=E_{d,144},\quad E_{0,0}=6000
\end{aligned}
\end{equation}

\textbf{步骤2：波动电价下的滚动调整模型}

对应问题三，0:00 计划购电的目标函数与功率平衡约束为
\begin{equation}
\begin{aligned}
&\min \; Z^{\mathrm{p}}=\sum_{k=0}^{143} g^{\mathrm{p}}_k\,p_{d,k}\,\Delta t\\
&g^{\mathrm{p}}_k+S^{0}_k+d^{\mathrm{p}}_k-c^{\mathrm{p}}_k\geq L_k
\end{aligned}
\end{equation}
调整子问题的分段费用上境约束、紧急购电缺口及其费用的核算与问题三完全相同，仅需将式~\eqref{eq:q3_Ck}、式~\eqref{eq:q3_w1}~和式~\eqref{eq:q3_w2}~中的电价 $p_k$ 逐项替换为 $p_{d,k}$。

\textbf{步骤3：全年总费用}

波动电价下，问题二对应的全年总购电费用与问题三对应的全年总费用分别为
\begin{equation}
\begin{aligned}
&Z^{(2)}=\sum_{d=0}^{364}Z_d\\
&Z^{(3)}=\sum_{d=0}^{364}\sum_{k=0}^{143}\Big(C_{d,k}+5\,p_{d,k}\,\delta_{d,k}\Delta t\Big)
\end{aligned}
\end{equation}
其中，$C_{d,k}$ 为式~\eqref{eq:q3_Ck}~所定义的第 $d$ 天第 $k$ 个间隔分段购电费用。该模型保持线性规划结构不变，电价参数 $p_{d,k}$ 的逐间隔替换即可反映实时波动电价对微网购电策略的影响。
